\chapter{Background}

\begin{figure*}[t]
\begin{pchstack}[boxed,center,space=1em]

\begin{pchstack}
  \procedure{$\main$ $\mathcal{G}^{\dl}_{\A}(\secp)$}{
  (\Gr, p , g) \randpick \GroupGen(\usecp) \\
  x \randpick \Zq;~X \gets g^{x} \\
  x' \randpick \A((\Gr, p , g), X) \\
	\pcif x' = x \\
	\pcind \pcreturn 1\\
	\pcreturn 0
  }
\end{pchstack}

\begin{pchstack}
  \vline
      \pchspace
  \procedure{$\main$ $\G^{\ell\hbox{-}\colorbox{gray}{$\mathsf{a}$}\omdl}_{\A}(\secp)$ }{
  (\Gr, p , g) \randpick \GroupGen(\usecp) \\
  Q \gets \emptyset \\
  q \gets 0 \\
  \pcfor i \in [0..\ell] \pcdo \\
  \pcind x_i \randpick \Zq;~X_i \gets g^{x_i} \\
  \pcind Q[X_i] = x_i \\
  \vec{x} \gets (x_0, \ldots, x_\ell) \\
  \vec{X} \gets (X_0, X_1, \dots, X_{\ell}) \\
  \vec{x'} \gets \A^{\BigO^{\dl}}((\Gr, p, g), \vec{X})  \\
  	\pcif \vec{x'} = \vec{x}~\wedge~q < \ell \\
	\pcind \pcreturn 1\\
	\pcreturn 0
  }
  \pchspace
  \procedure{$\BigO^{\dl}(X, \colorbox{gray}{$\alpha, \{\beta_i\}_{i=0}^{\ell}$} )$\  }{
    \pccomment{\colorbox{gray}{$X = g^{\alpha} \prod_{i=0}^{\ell} X_i^{\beta_i}$}} \\
  q \gets q + 1 \\
  %\colorbox{light-gray}{$\pcreturn \bot\ \pcif Q[X] = \bot$} \\
  \colorbox{gray}{$x \gets \alpha + \sum_{i=0}^{\ell} x_i \beta_i$} \\
  \colorbox{gray}{$\pcreturn x$} \\
  x \gets \mathsf{dlog}(X) \\
  \pcreturn x
  }
\end{pchstack}

\end{pchstack}
    \caption{The Discrete Logarithm (DL), One-More Discrete Logarithm (OMDL), and Algebraic One-More Discrete Logarithm (AOMDL) games.
    $\Gr$ is a cyclic group with prime order $q$ and generator $g$.
    $\mathsf{dlog}$ is an algorithm that finds the discrete logarithm relation of the input $X$ and $g$.
    The differences between the OMDL and AOMDL games are highlighted in gray;
    the key distinction is that in the AOMDL game,
    the adversary can only query for linear combinations of its challenge group elements;
    in the OMDL game, the environment must return the discrete logarithm of an
    arbitrary group element.}
    \label{fig:dl-omdl}
\end{figure*}


\section{Notation} \label{notation}

We represent a security parameter as $\lambda \in \mathbb{Z}$.
We denote the assignment of an element $y$ to the value $x$ as $y \gets x$,
and sampling an element from some set $S$ uniformly at random as $x \rgets S$.
For a randomized algorithm $\mathit{A}$, we write $x \rgets \mathit{A}()$ to indicate the random variable $x$ that is output from the execution of $\mathit{A}$.
We say that $f$ is a negligible function of $\lambda$ if it holds that
$\lim_{\lambda \to \infty} f(\lambda) \lambda^c = 0$,
for all $ c > 0$~\cite{BonehS2020}.

Let $\GG$ be a cyclic group of prime order $q$, and $\Zq$ be the field of integers modulo $q$.
Let $g$ be a generator of $\GG$.
We use $\set{n}$ to represent the set $\{1,\ldots,n \}$,
and represent integer vectors in bold-face font.

\paragraph{Polynomial Interpolation}
A polynomial $f(x) = a_0 + a_1 x + a_2 x^2 + \ldots + a_{\tcor} x^{\tcor}$
of degree $\tcor$ over a field $\F$ can be interpolated by $\thresh$ points.
Let $\eta \subseteq \set{n}$ be the list of $\thresh$ distinct indices
corresponding to the $x$-coordinates $x_i \in \F, i \in \eta$, of these points.
Then the Lagrange polynomial $L_i(x)$ has the form:

\begin{equation}\label{eqn:lagrange}
L_i(x) = \prod_{j \in \eta; j \neq i } \frac{x-x_j}{x_i - x_j}
\end{equation}

Given a set of $\thresh$ points $(x_i, f(x_i))_{i \in [\thresh]}$, any point $f(x_\ell)$ on the polynomial $f$ can be determined by Lagrange interpolation as follows:

\[ f(x_\ell) = \sum_{k \in \eta} f(x_k) \cdot L_k(x_\ell) \]

\paragraph{Game-Based Security.}
An experiment (or game) $\textsf{Exp}_{\adv, \mathcal{C}}^{\mathsf{sec}}$ is
played with respect to a security notion $\mathsf{sec}$,
an adversary $\adv$ modeled by a probabilistic (randomized) algorithm,
and a cryptographic primitive $\mathcal{C}$.
The experiment is initialized within a main procedure,
whose output is the output of the game.
The goal of $\adv$ is to win the game,
the conditions under which is defined within the experiment.
The probability of $\adv$'s success thus defines how secure $\mathcal{C}$ is with respect to the security notion $\mathsf{sec}$.

\paragraph{PPT Adversaries.}
We say that an adversary $\adv$ is probabilistic polynomial-time (PPT) if it is run on
random coins and its execution is performed in time that is polynomial relative to the size of its input.
In other words,
for every input $x \in \zo^n$,
$\adv(x)$ terminates within at most $f(n)$ steps,
for some polynomial $f$.

\section{Assumptions} \label{assumptions}

\paragraph{Discrete Logarithm (DL) Assumption.}
I show the DL assumption along with the DL assumption in Figure~\ref{fig:dl-omdl},
but give an intuition here.
To begin,
a PPT adversary is given a value $X \in \GG$ such that $X=g^x$ without knowing the value of $x$.
The challenge to the adversary is to output $x$.

We say that the DL assumption holds if the advantage to a PPT adversary in
successfully producing such an $x$ is a negligible function of $\lambda$.

\paragraph{Computational Diffie-Hellman Assumption (CDH).}
The Computational Diffie-Hellman Assumption (CDH) can be described as follows.
A PPT adversary is given the following values:

\[
  (g, g^a, g^b)
\]

The challenge to the adversary is then to compute $g^{ab}$.
The CDH assumption holds if the advantage to any PPT adversary in successfully
doing so is a negligible function of $\lambda$.

\paragraph{One-More Discrete Logarithm Assumption (OMDL).}
The One-More Discrete Logarithm (OMDL) assumption~\cite{BellareNPS03,BauerFP21} extends
the DL assumption, which is non-interactive, into an interactive assumption.
I show the OMDL assumption along with the DL assumption in Figure~\ref{fig:dl-omdl}.

In particular,
a PPT adversary is given $\ell+1$ discrete logarithm challenges $(X_0, \ldots, X_\ell) \in \GG^{\ell+1}$
with respect to a generator $g \in \GG$.
It is also given access to a discrete logarithm oracle,
that when queried on an arbitrary group element $X' \in \GG$,
returns the discrete logarithm of that value.
The adversary is said to win the OMDL game if it makes $\ell$ or fewer queries
to the discrete logarithm oracle,
and produces correct solutions $(x_0, \ldots, x_\ell)$ to its $\ell+1$ challenges such that $X_i = g^{x_i}$,
$i \in \{ 0, \ldots, \ell\}$.

The OMDL assumption holds if the advantage to any PPT adversary in winning the
OMDL game is a negligible function of $\lambda$.

\paragraph{Algebraic One-More Discrete Logarithm Assumption (AOMDL).}
The algebraic one-more discrete logarithm assumption (AOMDL) was formalized by Jonas, Ruffing, and Seurin~\cite{NickRS21}.
We highlight differences with the OMDL assumption in Figure~\ref{fig:dl-omdl}.
Note that while the OMDL assumption is not falsifiable
(because the adversary is allowed to query for the discrete logarithm of any arbitrary group element).
the AOMDL assumption is falsifiable.
This is because the input to the discrete logarithm oracle is a linear combination of all challenges issued to the adversary,
so it is guaranteed that the environment runs in polynomial time with respect to a PPT adversary.
Hence, the AOMDL assumption is a strictly weaker assumption than standard OMDL.

An adversary in the AOMDL game is initialized in exactly the same manner as an OMDL adversary,
with the same winning condition.
The only distinction between the games is the inputs to the discrete logarithm oracle.
In the AOMDL setting, the adversary provides a vector $(\alpha, \beta_0 \ldots, \beta_\ell)$ of coefficients to the discrete logarithm oracle.
The oracle then responds with the integer linear combination of these coefficients with respect to the
discrete logarithm of the set of challenges $(X_0, \ldots, X_\ell)$.
Informally, an AOMDL adversary can be used to win the OMDL game,
therefore implying that if OMDL is hard, then AOMDL is hard.

The AOMDL assumption holds if the advantage to any PPT adversary in winning the
AOMDL game is a negligible function of $\lambda$.

\paragraph{Algebraic Group Model (AGM).}
Introduced by Fuchsbauer, Kiltz, and Loss~\cite{FuchsbauerKL18},
the Algebraic Group Model (AGM) gives an idealized representation of a group,
by requiring the adversary to act algebraically.
In other words,
when an adversary outputs a group element $X \in \GG$ to the environment,
it additionally is required to output its representation with respect to all
other group elements the adversary has seen thus far.
In other words, the adversary must output a coefficient vector $\textbf{v} = (\alpha, b_1, b_2, \ldots)$,
such that $X = g^\alpha A_1^{b_1} A_2^{b_2} \ldots$.


\paragraph{Knowledge of Exponent Assumption (KoE).}

The Knowledge of Exponent (KoE) assumption~\cite{Damgard91,BellareP04} is an unfalsifiable assumption.
In short,
the assumption says that for an adversary that is given input $(g, h=g^s, q)$,
if the adversary outputs a value $(C, Y) \in \GG^2$ such that $C^s = Y$,
then there exists an extractor algorithm that given $(g, h=g^s, q)$,
%the adversary's state and $(C, Y) \in \GG^2$,
the extractor can output $x$ such that $C = g^x$ (and, consequently $Y = C^x$).
Intuitively, KoE assumes that in order to output such a $(C, Y)$,
the adversary must know $x$.


\section{Digital Signatures and Schnorr Signatures} \label{schnorr-sigs}

\subsection{Digital Signatures}

\begin{definition}[Digitial Signatures] \label{defn:ds}
A digitial signature scheme consists of polynomial-time algorithms $(\Setup, \KeyGen, \Sign, \Verify)$, defined as follows:

  \begin{itemize}

    \item $\Setup(1^\secp) \rightarrow \pp$:
    On input a security parameter $1^\secp$, this algorithm outputs public parameters $\pp$ (which are given implicitly as input to all other algorithms).

  \item $\KeyGen() \rightarrow (\pk, \sk)$:
  This probabilistic algorithm outputs a public key $\pk$ and secret key $\sk$.

  \item $\Sign(\sk, m) \rightarrow \sigma$:
    On input a secret key $\sk$ and a message $m$, this algorithm outputs a signature $\sigma$.

  \item $\Verify(\pk, m, \sigma) \rightarrow 0/1$:
  On input a public key $\pk$,
    a message $m$, and a purported signature $\sigma$, this deterministic algorithm outputs $1$ (accept) if $\sigma$ is a valid signature on $m$ under $\pk$; else, it outputs $0$ (reject).

\end{itemize}
\end{definition}

% standard EUF-CMA
\begin{figure}[t]
	\centering
	%\begin{linedgamespec}
	\begin{pchstack}[boxed,space=1em]
		\begin{pcvstack}[space=1em]
			\procedure{$\main$ $\G^{\eufcma}_{\A,\dsig}(\kappa)$}{
				\pp \leftarrow \ParGen(1^\kappa) \\
				Q_m \gets \emptyset \\
				(\pk, \sk) \randpick \KeyGen()\\
				(m^*, \sigma^*) \randpick \A^{\BigO^{\Sign}}(\pk)\\
				\pcreturn 0~\pcif m^* \in Q_m \\
				\pcind \vee \Verify(\pk, m^*, \sigma^*) \neq 1 \\
				\pcreturn 1
			}
		\end{pcvstack}
		\begin{pcvstack}[space=1em]
			\procedure{$\BigO^{\Sign}(m)$}{
				Q_m \gets Q_m \cup \{m\} \\
				\sigma \gets \Sign(\sk, m) \\
				\pcreturn \sigma
			}
		\end{pcvstack}
	\end{pchstack}
	%\end{linedgamespec}
	\caption{The EUF-CMA security game for a digital signature scheme $\dsig$.  The public parameters $\pp$ are implicitly given as input to all algorithms.}
	\hrulefill
	\label{fig:euf-cma}
\end{figure}

\noindent A digital signature scheme must be \emph{correct}.  It is said to be \emph{secure} if it is existentially unforgeable under
chosen-message attacks (EUF-CMA).

\begin{definition}[Correctness]
A digital signature scheme $(\Setup, \KeyGen, \Sign, \allowbreak \Verify)$ is \emph{correct} if for all security parameters $\secp \in \mathbb{N}$, all key pairs $(\pk, \sk) \in [\KeyGen()]$, and all messages $m$, we have:
\[ \Pr[ \Verify(\pk, m, \Sign(\sk, m)) = 1 ] = 1 \]
\end{definition}

\begin{definition}[EUF-CMA]\label{defn:euf-cma}
Let the advantage of an adversary $\A$ playing the Existential Unforgeability under Chosen-Message Attack (EUF-CMA) game,
$\G^{\suf}$, as defined in Figure~\ref{fig:euf-cma}, be as follows:
%
\[
  \advant{$\suf$}{\A,\dsig}(\secp) = \bigabs{\Pr[\G^{\suf}_{\A,\dsig}(\secp) = 1]}
\]
A digital signature scheme $\dsig = (\Setup, \KeyGen, \Sign, \allowbreak \Verify)$ is \emph{existentially unforgeable under chosen-message attacks} if for all PPT adversaries $\A$, $\advant{$\suf$}{\A,\dsig}(\secp)$ is negligible.
\end{definition}

\subsection{Schnorr Signatures}

A Schnorr signature is a Sigma protocol zero-knowledge proof of knowledge of the discrete logarithm of the public key $\pk$, made non-interactive and bound to the message $m$ by the Fiat-Shamir transform~\cite{FiatS86}.
Schnorr signatures are secure under the discrete logarithm assumption in the random oracle model~\cite{PointchevalS00}.


\begin{definition}[Schnorr Signatures~\cite{Schnorr91}] \label{defn:schnorr}
The Schnorr signature scheme consists of polynomial-time algorithms $(\Setup, \KeyGen, \Sign, \Verify)$, defined as follows:

  \begin{itemize}

    \item $\Setup(1^\secp) \rightarrow \pp$:
    On input a security parameter $1^\secp$, run $(\Gr, p, g) \gets \GroupGen(\usecp)$ and select a hash function $\Hash: \{0,1\}^* \rightarrow \Z_p$. Output public parameters $\pp \gets ((\Gr, p, g), \Hash)$ (which are given implicitly as input to all other algorithms).

  \item $\KeyGen() \rightarrow (\pk, \sk)$:
  Sample a secret key $\sk \randpick \Z_p$ and compute the public key as $\pk \gets g^x$.
    Output key pair $(\pk, \sk) \gets (\pk, \sk)$.

  \item $\Sign(\sk, m) \rightarrow \sigma$:
    On input a secret key $\sk = x$ and a message $m$, sample a nonce $r \randpick \Z_p$. Then, compute a
    nonce commitment $R \gets g^r$, the challenge $c \gets \Hash(\pk, m, R)$, and the
    response $z \gets r + cx$. Output a signature $\sigma \gets (R,z)$.

  \item $\Verify(\pk, m, \sigma) \rightarrow 0/1$:
  On input a public key $\pk = \pk$,
    a message $m$, and a purported signature $\sigma = (R, z)$, compute $c
    \leftarrow \Hash(\pk, m, R)$ and output $1$ (accept) if $R \cdot \pk^c = g^z$; else, output $0$ (reject).

\end{itemize}
\end{definition}


\section{Dodis-Yampolskiy VRF} \label{dy-vrf}

\section{Non-Interactive Key Exchange (NIKE)} \label{nike}

\section{Threshold Secret Sharing Schemes} \label{tss}

\begin{definition}{Threshold Secret Sharing}
A $(n, \thresh)$-threshold secret sharing scheme $\tss$ is the tuple of
  algorithms   $\tss=(\Setup, \IssueShares, \Recover)$, such that:
  \begin{itemize}
    \item  $\Setup(1^\secp) \rightarrow \pp$: On input the security parameter,
      output public parameters $\pp$ which are given implicitly as input to all other algorithms.

    \item $\IssueShares(\secp, s, n, \thresh) \rightarrow (\{ (1, \share_1), \ldots, (n, \share_n)\})$:
      A probabilistic algorithm performed that accepts as input
      the security parameter $\secp$,
      the secret $s$,
      the number of participants $n$, and the threshold $\thresh$,
      where $n$ and $t$ are positive integers, and $n \geq \thresh$.
      Outputs the list of shares $\{ (1, \share_1), \ldots, (n, \share_n) \}$.

    \item $\Recover(\thresh, \{ (i, \share_i)\}_{ i \in \coalition}) \rightarrow \sk/\fail$:
      A deterministic algorithm that accepts a threshold $\thresh$ and a set of shares $\{(i, \share_i) \}_{i \in \coalition}$,
      output $\fail$ if $\coalition \not\subseteq \set{n}$ or if $\lvert \coalition \rvert < \thresh$.
      Otherwise, recover $\sk$ using the set of shares.

    \item For correctness, for all security parameters $\secpar$ and $n, \thresh, \coalition$,
      such that $n \geq \thresh$,
      $\thresh \geq 1$,
      $\coalition \subseteq \set{n}$,
      and $\lvert \coalition \rvert \geq \thresh$,
      the following must hold:
    \begin{align*}
      & \tss.\Recover(\thresh, \{ (i, \share_i)\}_{ i \in \coalition}) = s,
      \text{ where }   \\
      & \tss.\IssueShares(\secpar, s, n, \thresh) \rightarrow ( \{ (i, \share_i) \}_{i \in \set{n}})
    \end{align*}
\end{itemize}
\end{definition}

\begin{remark}[Required number of shares for recovery]
Because $\recover$ requires a set of size at least $t$,
the definition implicitly ensures that the algorithm receives a sufficient
number of shares for recovery.
\end{remark}

\subsubsection{Shamir Secret Sharing}
We now define Shamir secret sharing as one concrete instantiation of a threshold secret sharing scheme.

\begin{definition}[Shamir secret sharing~\cite{Shamir79}] \label{defn:shamir}
Shamir secret sharing $\shamir$ is concrete instance of a threshold secret sharing scheme that consists of the following algorithms:

\begin{itemize}
    \item  $\Setup(1^\secp) \rightarrow \pp$: On input the security parameter,
      run $(\Gr, p, g) \gets \GroupGen(1^\secp)$.
      Output public parameters $\pp = (\Gr, p, g)$ which are given implicitly as input to all other algorithms.

    \item $\shamir.\IssueShares(s, n, \thresh) \rightarrow \{ (1, \share_1), \ldots, (n, \share_n)\}$:
      On input a secret $s$,
      the number of participants $n$,
      and a threshold $\thresh$, perform the following.
      First, define a polynomial $f(Z) = s + a_1 + a_2 Z^2 + \cdots + a_{\tcor} Z^{\tcor}$ by sampling $\tcor$ coefficients at random $(a_1, \ldots, a_{\tcor}) \randpick \Z_p$.
      Then, set each participant's share $\share_i, i \in \set{n}$, to be the evaluation of $f(i)$:
    \begin{equation*} \label{eq:issueshares}
        \share_i \gets s + \sum_{j \in \set{\tcor}} a_j i^j
    \end{equation*}

    Output $\{(i, \share_i) \}_{i \in \set{n}}$.

  \item $\shamir.\Recover(\thresh, \{ (i, \share_i)\}_{ i \in \coalition}) \rightarrow \sk/\fail$:
    On input a threshold $\thresh$ and a set of shares $\{(i, \share_i) \}_{i \in \coalition}$,
    output $\fail$ if $\coalition \not\subseteq \set{n}$ or if $\lvert \coalition \rvert < \thresh$.
    Otherwise, recover $s$ as follows:
    \begin{equation*} \label{eq:recovershares}
      s \gets \sum_{i \in \coalition} \lambda_i \share_i
    \end{equation*}
    where the Lagrange coefficient for participant $i$ in the set $\coalition$ is defined by
\[ \lambda_i = \prod_{j \in \coalition, j\neq i} \frac{j}{i - j} \]
\end{itemize}
\end{definition}

\section{Verifiable Secret Sharing} \label{verifiable-tss}

\emph{Verifiable Secret Sharing} (VSS) allows a dealer to share a secret $s$ in such a way that
participants can ensure that combining their shares allows for recovery of a secret that corresponds to a public commitment to $s$, without learning $s$ in the process.
%
%As such, VSS schemes are akin to public-key schemes, and rely on distinct secret and public domains,
%where both $s$ and participant shares are in the secret domain, and commitments to these shares and $s$ are in the public domain.
Participants verify that their shares can correctly recover $s$ using this commitment.
While prior definitions of VSS schemes~\cite{KatzLindell2014} define algorithms to issue and verify shares, and recover the secret,
we additionally define an algorithm to derive shares in a public domain.
This notion has been employed implicitly in prior DKGs~\cite{KomloG20}, and we formalize this construct here.

\begin{definition} \label{defn:vss}
  A \defn{Verifiable Secret Sharing (VSS) scheme} $\vess$ is the tuple
  $(\Setup, \issueshares, \recover, \verify, \derivepshare)$
  and parameterized by a one-way map $\vfunc : \secretspace \rightarrow \publicspace$ that maps elements in the secret domain $\secretspace$ to the public domain $\publicspace$,
  where
  \begin{itemize}
    \item $\Setup(1^\secp) \rightarrow \pp$:
    On input a security parameter $1^\secp$, this algorithm outputs public parameters $\pp$ (which are given implicitly as input to all other algorithms).
    \item $\IssueShares(\secret, n, t) \rightarrow (\{ (1, \share_1), \ldots, (n, \share_n)\}, \vsscommit)$:
      A probabilistic algorithm performed by a dealer that accepts as input
      the secret $\secret \in \secretspace$ from the secret domain $\secretspace$,
      the number of participants $n$, and the threshold $t$,
      where $n$ and $t$ are positive integers, and $n \geq t$.
      Outputs the list of shares $\{ (1, \share_1), \ldots, (n, \share_n) \}$
      and commitment $\vsscommit$.
    \item $\verify(i, \share_i, \vsscommit) \rightarrow \zo$:
      A deterministic algorithm performed by a recipient of a share that accepts
      an identifier $i$,
      a share $\share_i$,
      and a commitment $\vsscommit$.
      Outputs $1$ if the share is valid with respect to $\vsscommit$, otherwise outputs $0$.
    \item $\recover(t, \recoveryset) \rightarrow \secret/\fail$:
      A deterministic algorithm that accepts a recovery set $\recoveryset = \{ (j, \share_j)\}_{j \in C}$,
      where $C$ is a set of participant identifiers such that $\lvert C \rvert \geq t$.
      Employ $\recoveryset$ to obtain $s$,
      or output $\fail$ in the case of failure.
    \item $\derivepshare(i, \vsscommit) \rightarrow \pshare_i$:
      A deterministic algorithm that accepts an identifier $i$ and a commitment $\vsscommit$.
      Outputs a share in the public domain $\pshare_i$ associated with identifier $i$.
  \end{itemize}
\end{definition}

For correctness, for all security parameters $\lambda$,
for all $s \in \secretspace$, and $n, t, C$,
such that $n \geq t$, $t \geq 1$, $C \subseteq \set{n}$, and $\lvert C \rvert \geq t$,
the following must hold:
    \begin{align*}
      & \vess.\verify(i, \share_i, \vsscommit) = 1 \text{ for all } i \in \set{n},  \text{ and } \\
      & \vess.\derivepshare(i, \vsscommit) = \vess.\vfunc(\share_i) \text{ for all } i \in \set{n}, \text{ and}  \\
      & \vess.\recover(t, \recoveryset) = \secret,
      \text{ where }   \\
      & \vess.\issueshares(\lambda, \secret, n, t) \rightarrow ( \{ (i, \share_i) \}_{i \in \set{n}}, \vsscommit) \text{ and } \recoveryset = \{ (i, \share_i)\}_{i \in C}
    \end{align*}

\subsubsection{Feldman's VSS}
\label{feldman-vss}

As a concrete VSS in the discrete logarithm setting,
Feldman's Verifiable Secret Sharing~\cite{Feldman87} is an important building block.
Feldman's VSS itself is an augmentation to Shamir's secret
sharing~\cite{Shamir79}, allowing participants to verify the correctness of
their in zero knowledge, by committing to the polynomial ``in the exponent.''
Here, the map $\vess.\vfunc$ is simply $a \mapsto g^a$,
where $g$ is the generator of the group $\GG$, and
$a \in \Zq$.

We now define the scheme more precisely.

\begin{itemize}
    \item $\IssueShares(\lambda, s, n, t) \rightarrow (\{ (1,\share_1), \ldots, (n,\share_n)\}, \vsscommit)$.
      Accepts as input the security parameter $\lambda$,
      a secret $s \in \Zq$,
      the number of participants $n \in \mathbb{N}$,
      and the threshold $t \in \mathbb{N}$ such that $1 \leq t \leq n$.
      Perform the following steps:
    \begin{enumerate}
        \item  Sample $t-1$ coefficients at random: $(a_1, \ldots, a_{t-1}) \rgets \Zq^{t-1}$.
        \item Using $s$ as the constant term and $a_1, \ldots, a_{t-1}$ as the remaining coefficients, define the $(t-1)$-degree polynomial $f(x) = s + \sum_{i=1}^{t-1} a_i x^i$.
        \item Generate $n$ shares $\share_j \in \Zq $ by deriving $ \share_j \gets f(j),\ j \in \set{n} $.
        \item Define $A_0 = g^s$ and $A_i = g^{a_i}$, for $i \in \set{t-1}$.
        \item Define the VSS commitment $\vsscommit$ as the vector of commitments to the coefficients of $f$: $\vsscommit \gets ( A_0, \ldots, A_{t-1} )$.
        \item Output $(\{(i, \share_i) \}_{i \in \set{n}}, \vsscommit)$, where each $i$ represents the identifier of the recipient of $\share_i$.
    \end{enumerate}

    \item $\verify(i, \share_i, \vsscommit) \rightarrow \zo$:
      Accepts as input a participant identifier $i \in \set{n}$,
      the share $\share_i$ that belongs to participant $i$,
      and a VSS commitment $\vsscommit$ that is a tuple of group elements, such that $\lvert \vsscommit \rvert = t$.
      Verifies that $(i,\share_i)$ is a valid point on $f$ as follows:
    \begin{enumerate}
        \item Parse $\langle A_0, \ldots, A_{t-1} \rangle \gets \vsscommit$.
        \item Output 1 if Equation~\ref{eq:vss-verify} holds; otherwise, output $0$.
          \begin{equation} \label{eq:vss-verify}
            g^{\share_i} \iseq \prod_{k=0}^{t-1} A_k^{i^k}
          \end{equation}
    \end{enumerate}

  \item $\recover(t, \recoveryset) \rightarrow s/\fail$:
    Accepts as input the threshold $t$,
    a recovery set $\recoveryset = \{ (j, \share_j) \}_{j \in C}$ consisting of secret shares for a coalition $C \subseteq \set{n}$
    where $\lvert C \rvert \geq t$.
    If $\lvert C \rvert < t$,
    output $\fail$.
    Using the inputs, perform the following steps:
    \begin{enumerate}
      \item Derive Lagrange coefficients for each identifier $\{ L_j(0)\}_{j \in C}$.
        %\ian{Even if $|K| > t$?  What if the shares are inconsistent?}
        % CK: I added verification as a step, which prohibits this case. Same in the definition.
        \item Derive $\secret = \sum_{j \in C} \share_j \cdot L_j(0)$.
        \item Output $\secret$.
    \end{enumerate}

    \item $\derivepshare(i, \vsscommit) \rightarrow \pshare_i$:
      Accepts as input a participant identifier $i$ and a VSS commitment.
      Outputs $\pshare_i = g^{f(i)}$ by computing
      $\pshare_i \gets \prod_{k=0}^{|\vsscommit|-1} A_k^{i^k}$.
\end{itemize}



\section{Pseudorandom Secret Sharing} \label{pseudorandom-tss}
